\chapter{Ход работы}
\label{ch:chap2}

\section{Введение}

В предыдущей лабораторной работе я построил в MATLAB Simulink систему массового обслуживания M/M/1 с персональной очередью. СМО имела один генератор 
заявок, три сервер и одну очередь. В этой работе увеличим количество очередй до трех, чтобы у каждого сервера была своя очередь.

\section{Модификация предыдущей схемы}

Модифицируем СМО из лабораторной работы №11 следующим образом:

\begin{enumerate}
    \item Добавить 2 дополнительных очереди.
    \item Поставить на выходе из системы мультиплексор, который бы мультиплексировал выходы трех серверов.
    \item Добавить на выходе генератора демультиплексор, который бы "раскидывал" заявки по очередям.
    \item На выходе каждой очереди поставить мультиплексор на 2 входа, чтобы мультиплексировать выход демультиплексора и обратную связь.
\end{enumerate}

\section{Результат}

В результате модификации получена следующая схема:

\begin{figure}[H]
    \centering
    \includegraphics[width=1.0\textwidth]{smo.png}
    \caption{СМО с персональной очередью}
\end{figure}

\section{Графики}

\begin{figure}[H]
    \centering
    \includegraphics[width=1.0\textwidth]{plot1.png}
    \caption{Графики}
\end{figure}

На графике в левом верхнем угле - график сгенерированных заявок. Остальные 3 графика - кол-во вышедших из очередей заявок. Можно заметить, что нагрузка распределена равномерно. 
Также можно обратить внимание, что сумма вышедших из очередей заявок больше, чем
число сгенерированных заявок. Это связано с тем, что сервера отбрасывают заявки с вероятностью 0.5.

\endinput
